\documentclass[11pt]{article}
\usepackage[margin=1in]{geometry}
\usepackage{amsmath,amssymb,amsthm,mathtools,microtype}
\usepackage[hidelinks]{hyperref}
\hypersetup{
 pdftitle={Talagrand's ultimate matching conjecture for finite exponents},
 pdfauthor={Michael Goldman, Francesco Grotto, Dario Trevisan}}
\newenvironment{aiuse}
  {\renewcommand{\abstractname}{AI use}\begin{abstract}}
  {\end{abstract}}
\newtheorem{theorem}{Theorem}
\newtheorem{lemma}[theorem]{Lemma}
\newtheorem{proposition}[theorem]{Proposition}
\newcommand{\E}{\mathbb E}
\newcommand{\PP}{\mathbb P}
\newcommand{\R}{\mathbb R}
\newcommand{\Div}{\operatorname{div}}
\newcommand{\Tr}{\operatorname{tr}}
\newcommand{\diam}{\operatorname{diam}}
\newcommand{\dist}{\operatorname{dist}}
\newcommand{\ee}{\mathrm e}
\newcommand{\dd}{\,\mathrm d}
\title{Talagrand's ultimate matching conjecture\\for finite exponents}
\author{Michael Goldman \and Francesco Grotto \and Dario Trevisan}
\date{}
\begin{document}
\maketitle

\begin{aiuse}
This document was generated by an LLM through interaction with the authors,
who do not guarantee its correctness. It is made publicly available through
the GitHub repository~\cite{umc:repository}, together with a partially
revised version of the arguments, in the interest of transparency regarding
its development process. The authors plan to provide a fully digested
version of the arguments in the near future.
\end{aiuse}

\begin{abstract}
We prove the finite-exponent part of Talagrand's ultimate matching
conjecture for two independent uniform point clouds in the unit square.
For each finite critical pair \(1/r_1+1/r_2=1/2\), a single matching
has bounded coordinatewise exponential \(r_i\)-moments at scale
\(\sqrt{\log N/N}\) with probability tending to one.
We reveal the point counts through adaptive anisotropic partitions.
A local inverse of divergence controls the fresh current on each child,
and an exponential phase with a memory term controls the lag of the
routing metric. A terminal correction and an integer relay construction
produce a permutation of the original points.
\end{abstract}

\section{Statement and setup}

Let \(X_1,\dots,X_N\) and \(Y_1,\dots,Y_N\) be independent uniform
samples in the unit square. Write \(s_N=\sqrt{\log N/N}\).
The finite-exponent assertion of the ultimate matching conjecture
\cite{TalagrandProblems} is as follows.

\begin{theorem}\label{thm:finite}
Fix \(2\le r_1,r_2<\infty\) with \(1/r_1+1/r_2\ge1/2\).
There is a constant \(C=C(r_1,r_2)<\infty\) such that, with
probability tending to one as \(N\to\infty\), a single permutation
\(\pi\) satisfies
\[
 \frac1N\sum_{j=1}^N
 \exp\!\left[
   \left(\frac{|X_j^i-Y_{\pi(j)}^i|}{C s_N}\right)^{r_i}
 \right]\le2,\qquad i=1,2 .
\]
\end{theorem}

We prove the critical case \(1/r_1+1/r_2=1/2\), where both
exponents exceed two. The strictly subcritical case follows by
choosing a coordinatewise larger finite critical pair and reducing
the exponents. Indeed \(t^r\le 1+t^{r'}\) for \(r\le r'\),
and a fixed dilation absorbs the resulting constant in the modular.

\paragraph{Notation.}
For \(p\ge1\) the Luxemburg norm \(\|f\|_{\psi_p}\) is the
infimum of \(A>0\) for which the indicated probability space satisfies
\[
 \int\bigl(\exp((|f|/A)^p)-1\bigr)\le1.
\]
We use Lebesgue measure on the unit square, mass \(1/N\) on
matched pairs, and mass \(h^2\) per grid face of a given direction,
with zeros adjoined to make a probability space in the last case.
The triangle inequality for this norm follows from convexity of
the Young function. A current on a domain \(Q\) with divergence \(f\)
and zero normal flux satisfies
\(\int_Q J\cdot\nabla\varphi=-\int_Q f\varphi\)
for smooth tests on the closure, or equivalently after extending
both sides by zero off \(Q\).
We write \(A\preceq B\) for quadratic-form order and
\(a\lesssim b\) when \(a\le Cb\) with \(C\) independent of \(N\).

Choose a dyadic width \(h=2^{-m}\asymp s_N\) with sufficiently
large fixed lower comparison constant. Then
\(\mu=Nh^2\asymp m\asymp\log N\).
On the fixed \(h\)-grid, all red and blue counts belong to
\([c\mu,C\mu]\) with probability \(1-o(1)\), by the binomial
bound and a union bound. Denote averaging on this grid by \(P_h\).
Set \(s=h m^3\), which will be the smallest intrinsic side in
the adaptive subdivisions.

We use the following elementary conditional Bernstein estimate.
If \(\xi_j\) are conditionally independent and centered with
\(|\xi_j|\le b\) and \(\sum_j\E[\xi_j^2\mid\mathcal F]\le v\),
then
\[
 \log\E[\exp(u\textstyle\sum_j\xi_j)\mid\mathcal F]
 \le C u^2v\quad (|u|b\le1),
 \qquad
 \PP(|\textstyle\sum_j\xi_j|>t\mid\mathcal F)
 \le2\exp[-c\min(t^2/v,t/b)] .
\]
For the first bound apply \(e^z\le1+z+Cz^2\) for \(|z|\le1\)
to each summand and multiply the expectations. Optimizing the
exponential Markov bound over the permitted tilts gives the second.

\section{Finite subdivisions and a local divergence inverse}\label{sec:geometry}

The following geometric construction will be used in coordinate
charts obtained by determinant-one affine maps.
Suppose a parallelogram \(Q\) has diameter \(D\) and contains a ball
of radius \(c_0D\), with \(c_0>0\) fixed. Fix a small
\(\theta>0\). If \(D\le D_0s\) we leave \(Q\) terminal,
where \(D_0\) is a large constant.
Otherwise take the maximal squares from an \(s\)-dyadic lattice
which are contained in \(Q\), have side \(2s\le\ell\le\theta D\),
and have distance at least \(\gamma\ell\) from \(\partial Q\).
The distance constant \(\gamma\) is large and fixed.
Complete the partition by the intersections of uncovered
\(s\)-squares with \(Q\). A clipped \(s\)-square is assigned to
a full uncovered \(s\)-square within distance \(Cs\), and these
unions are terminal cells. For this assignment, move a point of
the clipped square toward the center of the interior ball by
a distance \(O(s)\) to find a ball of radius \(3s\) inside \(Q\).
Indeed convexity gives a ball of radius \(t c_0D\) at the
point at fraction \(t\) along this segment. There is a full
\(s\)-square in the smaller ball, within distance \(Cs\) of
\(\partial Q\). Choose \(\gamma\) larger than that constant.
The full square is uncovered because any admissible square
containing it has boundary distance at least \(2\gamma s\).
We take \(D_0\) large enough for the central ball
construction and for \(\theta D_0>2\). Each full square receives
only \(O(1)\) clipped pieces. Empty pieces are discarded, and
boundaries of area zero can be assigned by a fixed convention.

The nonterminal children are squares. Assign to each child its
side length \(\ell\), with \(\ell=s\) for a merged terminal cell.
In the coordinates of \(Q\), each child \(S\) satisfies
\begin{equation}\label{eq:whitney}
 |S|\asymp\ell_S^2,\qquad
 \diam S\lesssim\ell_S,\qquad
 \ell_S\asymp \max\{s,\dist(x,\partial Q)\}\quad(x\in S).
\end{equation}
The constants may depend on \(c_0,\theta,\gamma,D_0\).
For a maximal square, its next dyadic ancestor either meets the
boundary, fails the distance condition, or exceeds the upper
side cap. In the first two cases the ancestor's diameter bounds
the additional distance from any point in the child to the
boundary. In the third case the child's side is comparable to
\(D\). This proves the upper bound on distance in
\eqref{eq:whitney}, and the distance condition proves the lower
bound. The same argument for the \(2s\) ancestor of an uncovered
full \(s\)-square gives boundary distance \(O(s)\), since the
upper side cap cannot fail. The clipping and assignment preserve
this bound. Taking \(D_0\) larger if necessary makes all child diameters
at most \(C\theta D\). In particular all terminal children from
a split have intrinsic diameter \(O(s)\).
A cell which becomes terminal without a split has intrinsic
diameter \(O(s)\) in the chart in which termination is decided.

For a shape-regular convex parent \(Q\), take a smooth probability
density \(\eta_Q\) supported on a central ball of radius \(cD\).
Choose it by translating and dilating a fixed bump. For mean-zero
\(f\) on \(Q\) define its current by sending \(f(y)\dd y\) along
rays from \(y\) to \(z\), averaged against \(\eta_Q(z)\dd z\).
In a weak pairing this is
\[
 \int J_f\cdot\nabla\varphi
 =\int_Q f(y)\int_Q \eta_Q(z)
       [\varphi(z)-\varphi(y)]\,\dd z\,\dd y
 =-\int_Q f\varphi .
\]
Thus extension by zero gives exact divergence \(f\) and zero normal
flux. Its evaluation kernel is
\begin{equation}\label{eq:kernel}
 K(x,y)=(x-y)\int_1^\infty \eta_Q\bigl(y+t(x-y)\bigr)t\,\dd t .
\end{equation}
This formula follows by changing variables from the ray parameter.
Direct differentiation gives
\[
 |K(x,y)|\lesssim |x-y|^{-1},\qquad
 |\nabla_x K(x,y)|+|\nabla_y K(x,y)|\lesssim |x-y|^{-2}.
\]
On the support of the integrand \(t|x-y|\le D\).
The bump and its gradient are bounded by \(CD^{-2}\)
and \(CD^{-3}\), and either differentiation of the bump
costs at most a factor \(t\).
There is also the stronger support restriction
\begin{equation}\label{eq:support}
 K(x,y)=0\qquad\text{if}\qquad
 |x-y|>C\dist(x,\partial Q).
\end{equation}
Indeed, when \(z=y+t(x-y)\) is in the central ball,
\(x=(1-1/t)y+z/t\). A supporting line realizing the boundary
distance of \(x\) has distance at least \(cD\) from \(z\)
and nonnegative distance from \(y\). Therefore
\(\dist(x,\partial Q)\ge cD/t\), whereas \(|x-y|\le D/t\).

\begin{lemma}[Projected evaluations]\label{lem:projected}
Let \(P\) be averaging on the children of \(Q\). For \(x,x'\)
in the same child \(S\), of size \(\ell\), and a fixed \(0<\alpha<1\),
\begin{align*}
 \|P K(x,\cdot)\|_2 &\le C,\qquad
 \|P K(x,\cdot)\|_\infty\le C/\ell,\\
 \|P(K(x,\cdot)-K(x',\cdot))\|_2
 &\le C (|x-x'|/\ell)^\alpha,\\
 \|P(K(x,\cdot)-K(x',\cdot))\|_\infty
 &\le C\ell^{-1} (|x-x'|/\ell)^\alpha .
\end{align*}
The same bounds hold after subtracting the parent averages.
\end{lemma}
\begin{proof}
For either evaluation, \eqref{eq:support} and \eqref{eq:whitney}
put the support in \(B(x,C\ell)\). Every child meeting this
support has diameter \(O(\ell)\), by the upper distance comparison
and the Lipschitz property of boundary distance.
Any child meeting \(B(x,c\ell)\) has area at least \(c'\ell^2\).
For a nonterminal \(S\) this follows from the lower boundary
distance at \(x\), and for a terminal \(S\) from the minimum area.
Thus the projected support has area \(O(\ell^2)\). On its
cells near the singularity, the integral bound \(O(\ell)\)
and the area lower bound give an average of size \(O(\ell^{-1})\).
Farther away the kernel bound gives the same result.

For \(d=|x-x'|\le c\ell\), the integral of the kernel difference
on \(B(x,C\ell)\) is \(O(d[1+\log(\ell/d)])\).
Integrate the two singularities separately within distance \(2d\),
then use the derivative bound along the segment, which is in \(Q\).
Averaging on the near cells divides this bound by
\(\gtrsim\ell^2\). On the other cells the derivative bound gives
\(O(d/\ell^2)\) directly. The projected support area and
\(u(1+\log(1/u))\lesssim_\alpha u^\alpha\) prove the claims.
Larger \(d\) follows from the single-evaluation bounds.
Subtracting the parent average contracts \(L^2\) and at most
doubles \(L^\infty\).
\end{proof}

The same near and far integration, over the whole parent, gives
\begin{equation}\label{eq:lone}
 \int_Q |K(x,y)-K(x',y)|\,\dd y
 \lesssim d(1+\log(D/d)),\qquad d=|x-x'|\le D/2.
\end{equation}
Thus currents of bounded sources have continuous representatives
inside their parents. They extend continuously by zero on the
boundary, since \eqref{eq:support} gives
\(\int_Q|K(x,y)|\,\dd y\lesssim\dist(x,\partial Q)\).

\section{Phase and determinant-one selector}

Put \(a_i=1/r_i\), \(F_i(t)=(1+t^2)^{r_i/2}\) and
\[
 H(y)=F_1(y_1)+F_2(y_2),\qquad v=\nabla H,\qquad
 B(y)=\operatorname{diag}(H^{1-2a_1},H^{1-2a_2})+vv^{\mathsf T}.
\]
Use \(C(y)=\sqrt{\det B(y)}\,B(y)^{-1}\).
Criticality gives
\[
 \det B=H(1+W),\qquad W=\sum_i v_i^2 H^{-1+2a_i},
 \qquad 1+W\asymp H .
\]
For the last assertion, the upper bound follows from
\(|v_i|\lesssim H^{1-a_i}\). At large \(H\), at least one
\(F_i(y_i)\) is at least \(H/2\), and its squared derivative
is \(\gtrsim H^{2-2a_i}\). Bounded states are covered by the term
\(1\). Consequently
\begin{equation}\label{eq:metric}
 \det C=1,\qquad C_{ii}\lesssim H^{2a_i},\qquad
 C^{-1}=\chi\left[\operatorname{diag}(H^{-2a_i})+ww^{\mathsf T}\right],
 \quad w=\frac{v}{\sqrt H},\quad \chi=\sqrt{\frac{H}{1+W}}\asymp1 .
\end{equation}
For the diagonal estimate, \(C_{ii}=B_{jj}/\sqrt{\det B}\) with
\(j\ne i\), and \(B_{jj}\lesssim H^{2-2a_j}\). Then use
\(1-2a_j=2a_i\).

Let \(L(y,G)=\sqrt{\operatorname{Tr}(G C(y)^{-1})}\) for a
determinant-one positive definite matrix \(G\).
Under the local bounds \(G_{ii}\le M H(y)^{2a_i}\), the lag obeys
\begin{equation}\label{eq:laglip}
 |\partial_j L(y,G)|\lesssim_M H(y)^{1/2-a_j}.
\end{equation}
In fact \(L\) is the norm of the vector with entries
\(\sqrt\chi \sqrt{G_{ii}} H^{-a_i}\) and
\(\sqrt\chi G^{1/2}w\).
One has
\[
 |\partial_j H|\lesssim H^{1-a_j},\qquad
 |\partial_j W|\lesssim H^{1-a_j},\qquad
 |\partial_j \chi|\lesssim H^{-a_j},\qquad
 |\partial_j w_i|\lesssim H^{1/2-a_i-a_j}.
\]
The bounds for \(W\) and \(w\) follow by differentiating and using
\(F_i''(y_i)\lesssim H^{1-2a_i}\). The local bounds on the diagonals
give \(|G^{1/2}u|\le\sum_i\sqrt{G_{ii}}|u_i|\), proving
\eqref{eq:laglip}. The same estimate, with a changed constant,
holds along segments on which \(H\) is comparable to its value
at the initial point and where \(G\) satisfies the local bounds
at that initial point.

For a target \(C\) and old matrix \(G\), let \(G^+\) be the
following bounded metric update. Fix \(K>1\) and set
\[
 G^+=G^{1/2}\operatorname{clip}_{[K^{-1},K]}
                (G^{-1/2}CG^{-1/2})G^{1/2},
\]
where clipping acts on the eigenvalues.
The two relative eigenvalues are reciprocal. Hence
\(\det G^+=1\) and \(K^{-1}G\preceq G^+\preceq K G\).
The update also satisfies
\begin{equation}\label{eq:clip}
 G^+\preceq K^{-1}G+C.
\end{equation}
Indeed the scalar clipping function is bounded by \(K^{-1}+t\)
for \(t>0\). If the target is \(C(z)\), contraction of the
relative eigenvalue distance gives
\begin{equation}\label{eq:contraction}
 L(z,G^+)\le K^{-1/2}L(z,G)+C_K .
\end{equation}
To see this, if the old relative eigenvalues are \(e^d,e^{-d}\),
the new relative eigenvalues with respect to the target are
\(e^{(d-\log K)_+},e^{-(d-\log K)_+}\). Taking the square root of
their sum gives \eqref{eq:contraction}.

\section{Revealing empirical counts}\label{sec:reveal}

At stage \(k\), let \(\mathcal F_k\) be the history of all labeled
cell memberships revealed in both clouds. The next subdivision
is chosen using \(\mathcal F_k\). Use fixed measurable choices
for grid placements and ties. Conditional on \(\mathcal F_k\),
the points remain independent and uniform in their respective
cells. For the induction step, after a subdivision has been
chosen, conditioning an independent parent-uniform position on
its child label makes it uniform in that child.
Terminal cells keep their counts without further revelation.

Let \(f_k\) be the signed count density divided by \(\sqrt N\)
on the partition at stage \(k\). Thus on \(Q\) it is
\((N_Q^X-N_Q^Y)/(\sqrt N|Q|)\).
Write \(\delta=\epsilon^2/m\), with \(\epsilon>0\) fixed later.
We build \(U_k\) with divergence \(\sqrt\delta f_k\) and zero
normal flux. At the root \(U_0=0\) and \(G_0=I\).
Write \(G_k(x)\) for the metric of the cell containing \(x\) at stage
\(k\), and suppress the spatial argument when the cell is understood.
Each nonterminal parent \(Q\) is a square in its own
\(G_k^{-1/2}\) coordinates. At its center \(x_Q\) select the
target \(C(U_k(x_Q))\), and form \(G^+\) by the bounded update.
In the \(G^{+\,-1/2}\) chart the parent is a parallelogram of
uniform shape regularity, by the relative matrix bound.
Use the subdivision of Section~\ref{sec:geometry} in this chart and assign \(G^+\)
to all its children. If the parent is too small to split,
make it terminal and keep the old field and metric there.
For each split, route the fresh source \(f_{k+1}-f_k\)
inside the parent using \eqref{eq:kernel} in the new chart.
It has mean zero in the parent. The determinant-one change
back to physical coordinates multiplies the current by
\(G^{+\,1/2}\) and preserves the divergence under composition
with the coordinate map, by a change of variables in the
weak identity. Multiply the increment by \(\sqrt\delta\).
The resulting field \(U_{k+1}\) has the required divergence.

Choose the child side cap \(\theta D\) so small, in terms
of \(K\), that every nonterminal child side is at most
\(\theta_0\) times the parent side in their respective
charts, where \(\theta_0\sqrt K<1\). Every nonterminal
cell has intrinsic side at least \(s\). Hence each branch
terminates in at most \(J_0\le C_{\rm dep}m\) stages.
Pad to this deterministic horizon, doing nothing on
terminated branches. All cells have area at least \(c s^2\).
In its chart of formation each terminal child has diameter
\(O(s)\). A parent declared terminal without a split has
diameter \(O(s)\) in the proposed chart.

All local estimates are used under
\(N_Q^X,N_Q^Y\le2N|Q|\). A stronger uniform event is
available on the attempted stages. The conditional mean of
a new child count is \(N_Q |S|/|Q|\). The error threshold
\(m^{-2}N|S|\) fails with conditional probability at most
\(2\exp(-cN|S|m^{-4})\) under the parent cap, by Bernstein.
Since \(N|S|\gtrsim Ns^2\asymp m^7\), the union failure
over at most \(Cs^{-2}\) children per stage and \(O(m)\)
attempted stages is \(o(1)\).
Starting at the root, on the good count event the density
of each cloud differs from \(N\) by a relative \(O(m^{-1})\).
This follows by summing the relative-to-\(N\) error
\(m^{-2}\) along the branch. Count tests may be applied
and the construction stopped after each attempted reveal.

In a split parent, let \(V_{k+1}=U_{k+1}-U_k\) be the Piola
current, and let \(Z_{k+1}=(G^+)^{-1/2}V_{k+1}/\sqrt\delta\).
We use intrinsic coordinates for the estimates on \(Z_{k+1}\).
For each test point, \(Z_{k+1}\) is a centered sum of independent
point contributions. Its evaluation against a direction is the
sum of \(N^{-1/2}\) times the child-averaged kernel evaluated at
the point's label, centered by the parent average, with a sign
for the cloud. By Lemma~\ref{lem:projected} and the parent cap,
each directional evaluation has total conditional variance \(O(1)\)
and summand bound \(C/(\sqrt N\ell)\) in a child of size \(\ell\).
The corresponding evaluation difference at distance \(d\) has
variance \(O((d/\ell)^{2\alpha})\) and summand bound
\(C(d/\ell)^\alpha/(\sqrt N\ell)\). In particular
\begin{equation}\label{eq:mgf}
 \log\E[\exp(u\cdot Z_{k+1}(x))\mid\mathcal F_k]\le C|u|^2,
 \qquad |u|\le c\sqrt N s .
\end{equation}

\begin{lemma}[An envelope on each new child]\label{lem:envelope}
Fix \(0<\beta<\alpha<1\). In each child \(S\), there is an envelope
\(E_{k+1,S}\) for the supremum and the scaled \(\beta\)-H\"older
seminorm of \(Z_{k+1}\) in intrinsic coordinates. For each fixed
\(R\), uniformly in old data satisfying the count cap, it has
conditional tail
\[
 \PP(E_{k+1,S}>t\mid\mathcal F_k)\le C\ee^{-ct^2},
 \qquad 0\le t\le R\sqrt m ,
\]
for sufficiently large \(m\). The constants \(c,C\) can be fixed
before choosing \(R\).
\end{lemma}
\begin{proof}
Take nets of mesh \(2^{-j}\ell\) inside \(S\) with cardinality
\(O(2^{2j})\). Test pairs at distance \(O(2^{-j}\ell)\)
in each net, each point against a nearest predecessor,
and an anchored value. At most \(O(2^{4j})\) tests are
needed at level \(j\).
At level \(j\) use the threshold
\(D 2^{-\alpha j}(1+t+j)\) for differences.
Bernstein gives a union failure bounded by \(C\ee^{-ct^2}\)
over all levels, because the exponent before counting pairs
is bounded below by
\[
 c\min\bigl(D^2(1+t+j)^2,\ D\sqrt N s(1+t+j)\bigr)
 \ \ge\ c'D(1+t^2+j)
\]
when \(t\le R\sqrt m\) and \(m\) is large.
Here \(D\) is chosen sufficiently large. Summation of the net
increments gives the supremum bound. For two points at distance
comparable to \(2^{-j}\ell\), sum each chain down to that scale
and then use the same-level comparison. The bound on the sum is
\(C2^{-\beta j}(1+t)\), using \(\beta<\alpha\), which proves
the seminorm bound. The evaluation field is continuous by \eqref{eq:lone}
applied to the bounded realized density.
The construction works on terminal unions since only the
diameter bound and the evaluation estimates on the same child
are used. Rescaling the threshold by a fixed factor gives
the stated tail.
\end{proof}

For a fixed spatial \(x\), write \(E_i(x)\) for the envelope of
its child at stage \(i\), and use zero on an inactive step.
The child is chosen before the fresh labels, so
Lemma~\ref{lem:envelope} applies conditionally to \(E_i(x)\).
Taking \(R\) sufficiently large and using the bound on the
number of children gives
\begin{equation}\label{eq:globalmax}
 E_{i,S}\le R\sqrt m
 \quad\text{for all attempted stages and all their children}
\end{equation}
with probability \(1-o(1)\), up to the count cap.

\section{Local diagonal bootstrap}

We prove by induction on the good history, pointwise on every
cell, that
\begin{equation}\label{eq:caps}
 G_{k,ii}\le M H(U_k(x))^{2a_i},\qquad
 (1-\eta) H(U_{k-1}(x))\le H(U_k(x))
       \le (1+\eta)H(U_{k-1}(x)).
\end{equation}
The ratio in \eqref{eq:caps} is asserted after a step, including
on terminal children of that step. It is automatic on inactive
steps. Here \(M\) is a large fixed constant and \(\eta>0\) is
small and fixed. The root has \(U_0=0\) and \(G_0=I\).
The selection and partition are well-defined without the local
caps. The geometric estimates use the relative bound \(K\) only.

For an active parent \(Q\) at stage \(k\), the ancestor H\"older
bound gives, for \(x,z\in Q\),
\[
 |U_k^i(x)-U_k^i(z)|
 \le C\sqrt\delta \sum_{j=1}^k
       \rho^{k-j}\sqrt{G_{j,ii}(x)} E_j(x), \qquad \rho<1 .
\]
Indeed a descendant diameter in the chart of its ancestor child
at stage \(j\) is at most \(C(\theta_0\sqrt K)^{k-j}\) times the
side of that child. This uses side shrink and
\(G_k\preceq K^{k-j}G_j\). Raising this factor to the H\"older
exponent gives \(\rho^{k-j}\). For a newly formed terminal
child use its diameter bound, with a fixed constant at the last
step. For a frozen leaf use the stage of its last refinement.
By induction
\(H(U_j(x))\le (1-\eta)^{-(k-j)}H(U_k(x))\).
Choose \(\eta\) small enough that
\(\rho'=\rho(1-\eta)^{-\max_i a_i}<1\).
Writing \(H=H(U_k(x))\), we obtain
\begin{equation}\label{eq:osc}
 |U_k^i(x)-U_k^i(z)|
 \le C \sqrt{\delta M} H^{a_i}
           \sum_{j=1}^k(\rho')^{k-j} E_j(x)
 \le C \epsilon\sqrt M R\,H^{a_i}.
\end{equation}
For the last inequality use the good envelope history.
If \(|z_i-y_i|\le b H(y)^{a_i}\) for both coordinates and
\(b\le1\), then \(H\) is bounded by \(C H(y)\) on the segment
from \(y\) to \(z\), and
\[
 |H(z)-H(y)|\le C b H(y).
\]
The first assertion follows by the power bounds on \(F_i\) and
the second by integrating the derivatives on the segment.
In particular \eqref{eq:osc} makes the representative phase
\(H(U_k(x_Q))\) within relative error \(O(\epsilon\sqrt M R)\)
of \(H\).

By \eqref{eq:clip} and \eqref{eq:metric},
\[
 G^+_{ii}\le
 \bigl[K^{-1}M+C(1+C\epsilon\sqrt M R)^{2a_i}\bigr]H^{2a_i}.
\]
Fix \(q_d\in(K^{-1},1)\) and \(M\) large enough that this is at
most \(q_d M H^{2a_i}\) for sufficiently small \(\epsilon\).
Choose \(\eta\) also so that
\(q_d\le (1-\eta)^{2\max_i a_i}\).
On the fresh local envelope cutoff, the increment satisfies
\[
 |V_{k+1}^i(x)|\le \sqrt\delta\sqrt{G^+_{ii}}\,E_{k+1}(x)
 \le C\epsilon\sqrt M R\, H^{a_i}.
\]
The relative-change estimate above gives the new ratio in
\eqref{eq:caps}, by taking \(\epsilon\) small. The margin \(q_d\)
then gives \(G^+_{ii}\le M H(U_{k+1}(x))^{2a_i}\).
This proves the induction. Only the fresh local envelope cutoff
is required for the pointwise conclusion after an old good history.

\section{Drift of the anisotropic phase}

For a fixed \(x\) in an active parent, put
\[
 T_k=\sum_{j=1}^k \rho_1^{k-j} E_j(x)^2,\qquad
 L_k=L(U_k(x),G_k),\qquad S_k=L_k^2+T_k ,
 \qquad \rho'<\rho_1<1 .
\]
The sum in \eqref{eq:osc} is \(O(\sqrt{T_k})\) by Cauchy--Schwarz.
Use \eqref{eq:contraction} at the representative, and apply
\eqref{eq:laglip} to move between it, \(U_k(x)\) and \(U_{k+1}(x)\).
The local caps hold on these segments up to constants by the
relative smallness estimates in the preceding section. Thus on
the fresh local cutoff,
\begin{equation}\label{eq:lag}
 L_{k+1}\le K^{-1/2} L_k+C
       +C_M \epsilon \sqrt{H/m}\,(\sqrt{T_k}+E_{k+1}(x)).
\end{equation}
For example the coordinate error is bounded by
\(C\sqrt{\delta M} H^{a_i}(\sqrt{T_k}+E_{k+1})\), while the
lag derivative in that coordinate costs \(C_M H^{1/2-a_i}\).

Use the phase
\[
 \Phi_k=\lambda_k H(U_k)(1+\kappa\delta S_k),
 \qquad \lambda_k=\ee^{-\Lambda k\delta},
\]
with spatial modular overshoot threshold \(T_{\rm stop}\) fixed.
Freeze \(S_k\) on inactive steps. Before stopping, \(H\le C_H m\)
uniformly. Indeed \eqref{eq:osc} on the current leaf makes \(H\)
at every point of that leaf comparable to \(H\) at any given
point of it. A leaf has area \(\gtrsim s^2\) and \(\lambda_k\ge1/2\)
will be arranged. The modular bound gives \(H\le C_H m\), with
\(C_H\) fixed. For old frozen leaves use the estimate when they
were formed. This argument uses the inductive local caps and
does not require a preexisting global matrix cap.

By \eqref{eq:lag}, squaring with fixed slack and updating
\(T_{k+1}=\rho_1T_k+E_{k+1}^2\) gives
\begin{equation}\label{eq:memory}
 S_{k+1}\le q_0 S_k+C_0(1+E_{k+1}^2),\qquad q_0<1.
\end{equation}
One can choose \(q_0\) strictly between \(\max(K^{-1},\rho_1)\)
and 1, and then take \(\epsilon\) small. In the initial choices
take \(\rho_1\) fixed above the original geometric rate \(\rho\),
then take \(\eta\) small so that \(\rho'<\rho_1\) and
\((1+\eta)q_0<1\).
Since \(H'\le(1+\eta)H\) on the local cutoff,
\[
 H'S_{k+1}-HS_k\le-c H S_k+C H(1+E_{k+1}^2).
\]
There is no additional higher power from the memory product.

Write \(y=U_k(x)\) and \(V=V_{k+1}(x)\). On the local cutoff,
\(H\) stays comparable to its old value along the segment from
\(y\) to \(y+V\). Since each
\(F_i''\lesssim H^{1-2a_i}\) along that segment, Taylor gives
\[
 H(y+V)-H(y)\le v(y)\cdot V+C\sum_i H^{1-2a_i}|V_i|^2.
\]
Write \(\lambda=\lambda_{k+1}\), \(H=H(U_k)\),
\(E=E_{k+1}(x)\), \(Z=Z_{k+1}(x)\) and
\(\ell=\lambda\sqrt\delta (G^+)^{1/2}v(y)\).
Since \(vv^{\mathsf T}/H\preceq C_1 C(y)^{-1}\)
and \(G^+\preceq K G_k\),
\[
 |\ell|^2\lesssim \lambda^2\delta H L_k^2 .
\]
By \eqref{eq:caps}, \(L_k^2\lesssim_M H\), so the conditional
linear MGF at \(2\ell\) is within its allowed range
\(c\sqrt N s\). The Taylor quadratic term is bounded by
\(C_M\lambda\delta H E^2\).
It follows that on the local cutoff
\[
 \Phi_{k+1}-\Phi_k\le
   -(\lambda_k-\lambda)H-c\kappa\lambda\delta H S_k
   +C\kappa\lambda\delta H+\ell\cdot Z+uE^2,
 \qquad
 u\lesssim_{M,\kappa}\lambda\delta H .
\]
Since \(H\le C_Hm\), \(u\) is small for sufficiently small
\(\epsilon\). Cauchy--Schwarz separates the conditional linear
MGF from the truncated quadratic factor. The latter uses
\[
 \E[\ee^{2u E^2}\mathbf{1}_{E\le R\sqrt m}\mid\mathcal F_k]
 \le 1+\int_0^{R\sqrt m}4ut\ee^{2ut^2}
                 \PP(E>t\mid\mathcal F_k)\,\dd t
 \le 1+Cu
\]
for \(2u<c\) bounded away from \(c\). The combined log cost is
\(O(|\ell|^2+u)\), hence
\(O(\lambda\delta H S_k)+O(\lambda\delta H)\).
Choose \(\kappa\) large to pay the first term, then \(\Lambda\)
large to pay the second from \(\lambda_k-\lambda\ge
\Lambda\lambda\delta\). The conditional modular decreases on the
cutoff. The constants can be chosen in this order. Fix the geometry,
choose \(q_d>K^{-1}\), choose \(\rho_1>\rho\), and choose
\(q_0>\max(K^{-1},\rho_1)\), all below one. Choose \(\eta\)
for the three strict margins in the oscillation, cap and
memory estimates. Fix \(M\) large enough for the cap margin
assuming the representative phase ratio is at most two.
Then choose \(\kappa,\Lambda,T_{\rm stop}\).
Use \(\lambda_k\ge1/2\) provisionally for the bootstrap \(C_H\).
Choose \(R\) for the conditional union bound and finally take
\(\epsilon\) small enough for every smallness constraint, including
\(\Lambda J_0\epsilon^2/m\le\log 2\).
The relative oscillation estimate ensures the representative
ratio bound once \(\epsilon\) is small.

Define the scalar process by the spatial modular while active.
On count or envelope failure set it to zero and keep it there.
On modular overshoot keep the attained value thereafter.
At a point on a terminal branch, the phase decreases because
the field and memory freeze and the clock decreases.
Conditional integration of the pointwise bound proves the
supermartingale property because the simultaneous fresh good event
has indicator at most the local envelope cutoff used in the
pointwise estimate. The initial modular is bounded by a fixed
constant, since \(H(0)=2\), \(G_0=I\), and the initial lag is
fixed. A nonnegative supermartingale starting at \(B_0\)
crosses \(T_{\rm stop}\) with probability at most
\(B_0/T_{\rm stop}\), by stopping at first crossing.
Choose \(T_{\rm stop}\) large and fixed. Count and envelope
failures have probability \(o(1)\), so the probability of
completion with bounded modular is bounded below.
On this event \(\int\exp(\tfrac12|U_{J_0}^i|^{r_i})\le
T_{\rm stop}\). Jensen applied to a small power of this
exponential gives a fixed Luxemburg bound for \(U_{J_0}^i\).
The unnormalized leaf current \(J_{\rm leaf}=U_{J_0}/\sqrt\delta\)
therefore has the two coordinate norms \(O(\sqrt m)\).

\section{Correction to the fixed grid}\label{sec:correction}

On success the diagonal caps and \(H\lesssim m\) give
\(\|G_k\|\lesssim m^{2a_*}\), where \(a_*=\max_i a_i\).
This also bounds the proposed metric in which a parent is
declared terminal without a split. The physical terminal
leaf diameters are at most \(D_*=C s m^{a_*}\).
Let \(f_{\rm leaf}\) be their signed count density divided by
\(\sqrt N\), and let
\(f_{\rm emp}=N^{-1/2}(\sum_j\delta_{X_j}-\sum_j\delta_{Y_j})\).
The leaf current has exact divergence \(f_{\rm leaf}\).
To reach the fixed histogram \(P_h f_{\rm emp}\) we add
currents for
\[
 (P_h-I)f_{\rm leaf},\qquad P_h(f_{\rm emp}-f_{\rm leaf}).
\]
For the first source use the averaged-ray inverse on
each \(h\)-square, where the source has mean zero.
The \(L^1\) kernel bound yields current supremum at most
\(Ch\|f_{\rm leaf}\|_\infty\). The relative count event gives
\(\|f_{\rm leaf}\|_\infty\lesssim\sqrt N/m\), so this is
\(o(\sqrt m)\).

Use the averaged-ray inverse on the whole unit square for
the second source. Denote its evaluation kernel by \(K_x\),
a function of the source variable. The test against a
source atom uses \(P_hK_x\), which satisfies
\[
 |P_hK_x(z)|\lesssim (h+|x-z|)^{-1}.
\]
This follows by integrating the \(1/|x-z|\) singularity over
the cell when \(|x-z|=O(h)\), and the kernel bound away from it.
Condition on the full terminal revelation, including the
adaptive tree and success. Each point is independent and
uniform in its leaf. Its contribution to the correction
is centered by the leaf average of \(P_h K_x\), with magnitude
at most \(C/(h\sqrt N)\) after normalization. By the leaf
density cap, the total evaluation variance is bounded by
\[
 C\sum_{Q\ {\rm leaf}}\int_Q |P_hK_x-(P_hK_x)_Q|^2.
\]
For leaves meeting \(B(x,C(D_*+h))\), bound by the squared
regularized kernel on a slightly larger ball, at cost
\(C(1+\log_+(D_*/h))\). On the other leaves, the
source-variable derivative bound for \(K_x\) gives oscillation
\(O((D_*+h)/r^2)\) at distance \(r\) from \(x\).
Indeed each pair of grid cells meeting one leaf is at most
\(D_*+O(h)\) apart, so the same bound holds for their
averages. The squared integral over this far region is
\(O(1)\). Thus the variance is bounded by a deterministic
\(v=C(1+\log m)\).

For the correction current \(J_{\rm corr}\), Bernstein gives the
conditional pointwise tail
\[
 \PP(|J_{\rm corr}(x)|>t\mid\mathcal F_{J_0})
   \le C\exp\bigl[-c\min(t^2/v,h\sqrt N\,t)\bigr]
\]
on successful histories.
On the fixed-grid occupancy event its source density is
bounded by \(C\sqrt N\). The evaluation modulus \eqref{eq:lone}
then shows that a net of mesh \(N^{-2}\) controls the supremum,
up to an \(o(1)\) error. At threshold \(R_1\sqrt{mv}\) with
large fixed \(R_1\), Bernstein and a union bound give
\(\|J_{\rm corr}\|_\infty=O_{\PP}(\sqrt{mv})\) on old success.
Up to that cutoff the conditional pointwise tail is bounded
by \(C\exp(-ct^2/v)\) for large \(N\).
Tail integration gives a bounded expected truncated spatial
quadratic exponential modular at scale \(C\sqrt v\).
Markov's inequality and a fixed dilation of a bounded
modular give \(\|J_{\rm corr}\|_{\psi_2}=O_{\PP}(\sqrt v)\)
on the supremum event. If \(\|J\|_\infty\le A\) and
\(\|J\|_{\psi_2}\le B\), then
\(\|J\|_{\psi_r}\le A^{1-2/r}B^{2/r}\) for \(r\ge2\).
At this scale the \(r\)-th power is bounded pointwise
by \((|J|/B)^2\). It follows that for each fixed finite \(r_i\)
the correction has \(\psi_{r_i}\) norm
\(O_{\PP}(m^{1/2-1/r_i}\sqrt{\log m})=o_{\PP}(\sqrt m)\).
All corrections have exact divergence and zero normal flux.
Together with the leaf construction they give, with
probability bounded below, a current for \(P_h f_{\rm emp}\)
with both coordinate norms \(O(\sqrt m)\).

\section{From the current to a single point permutation}

Write \(\mu=Nh^2\), \(g=R-B\) for the raw count difference
on the fixed grid, and orient its faces. The grid divergence
is net outgoing load. We record the transfer and probability
upgrade so the proof ends on the original labeled points.

\subsection*{Extraction of a face flow}
In one dimension let \(I_0\) be linear interpolation through
the cell centers, constant on the two outer half-cells.
Write \(D_0\) for the unscaled forward difference on center
values. The map from center values to cell averages is
\(B_0=I-D_0^{\mathsf T}D_0/8\). The identity follows by
integrating each hat on the two half-cells, with the
constant extension at the ends. Put
\(B_1=I-D_0D_0^{\mathsf T}/8\), so \(D_0B_0=B_1D_0\).
Both inverse matrices have absolute row and column sums
at most 2 by the Neumann series.

For a grid test \(u\), apply \(I_0\) in both coordinates
to \((B_0^{-1}\otimes B_0^{-1})u\).
The resulting test has cell averages \(u\).
Its horizontal derivative is \(h^{-1}\) times the
interpolation of \((B_1^{-1}\otimes B_0^{-1})D_1u\),
constant between centers in the horizontal variable and
linear in the transverse variable. The vertical derivative
has the exchanged formula. Each product basis function for
the derivatives has mass \(h^2\), and the basis functions
sum to at most one. For a direction \(i\), write \(\beta_e\)
for these basis functions and \(T_i\) for the inverse tensor
filter. For a current \(J\) of divergence
\(P_h f_{\rm emp}=g/(\sqrt N h^2)\), put
\[
 A_{i,e}=h^{-2}\int J_i\beta_e,\qquad
 q_i=\sqrt N h\,T_i^{\mathsf T} A_i.
\]
The weak divergence equation yields
\(\sum_i\langle q_i,D_i u\rangle=-\langle g,u\rangle\).
For faces oriented in the forward direction this means
\(-\sum_i D_i^{\mathsf T}q_i=g\), the net outgoing divergence.
These piecewise bilinear tests are Lipschitz and may be used
in the weak identity by extending and mollifying. In the
Jensen estimate, dividing the absolute entries of a filter
by its maximal row sum gives subprobability weights.
The column sum bound controls the sum over faces.
Thus each coordinate face norm is at most
\(C\sqrt N h\) times the continuum coordinate norm.
Each face of a fixed direction carries mass \(h^2\) in this
Luxemburg norm, with zeros adjoined if needed.
Consequently the raw count load scale is
\(\sqrt N h\sqrt m\asymp\mu\).

\subsection*{Upgrade to high probability}
Let \(Q(g)\) be the minimum, over real flows of divergence \(g\),
of the maximum of the two coordinate face \(\psi_{r_i}\) norms.
This is a convex norm on mean-zero grid sources and is uniformly
Lipschitz in their Euclidean norm. In fact the elementary square
grid isoperimetric inequality
\(|\partial S|\gtrsim\sqrt{\min(|S|,|S^c|)}\) and Cauchy--Schwarz
give \(|\sum_{v\in S}g_v|\le C\|g\|_2|\partial S|\).
For the isoperimetric inequality, if neither a row nor a column
of \(S\) is full, all occupied rows and columns are mixed,
and their counts have product at least \(|S|\).
If some row is full and another is empty, every column is
mixed. If there is a full row and no empty row, the complement
has neither a full row nor a full column. The case of a full
column follows by exchanging coordinates. A flow with uniform capacity \(C\|g\|_2\) exists by these
cut inequalities. Indeed give each edge two directed arcs
with that capacity, attach a source to the positive
divergences and a sink to the negative ones, and maximize
the total sent. A maximizing flow exists by compactness.
If the demand has not been met, the vertices reachable
from the source in the residual network define a cut
of capacity below the demand. A residual augmenting path
would increase the flow, so none exists. This contradicts
the cut inequalities. Subtract opposite arc loads to
obtain the desired signed flow. The bound on its norm
and the triangle inequality prove the Lipschitz bound.

For a replacement of one point by an independent uniform
point, let \(g'\) be the new source and choose a supporting
vector \(u\) for \(Q\) at \(g\) in the mean-zero subspace.
Then \(\|u\|_2\le C\) and
\((Q(g)-Q(g'))_+^2\le 2(u_{\rm old}^2+u_{\rm new}^2)\)
by convexity. Sum over the \(2N\) replacements.
The old-cell terms are bounded by the maximum pooled occupancy
times \(C\). The fresh-cell terms cost \(O(Nh^2)\) after
conditional expectation. The expected maximum pooled occupancy is
\(O(\mu)\) since the exponential moment of the count in one cell is
at most \(\exp(C\mu)\) at tilt one. Jensen thus bounds the expected
maximum by \(\log(h^{-2})+C\mu\).
For independent inputs \(W_j\) and a function \(F\) of them,
the reveal martingale variance decomposition and conditional
Jensen give
\(\operatorname{Var}F\le\sum_j \E[\operatorname{Var}(F\mid W_{-j})]\),
where \(W_{-j}\) denotes all inputs except \(j\).
To see this, the \(j\)-th martingale difference is the
conditional expectation, given \(W_1,\ldots,W_j\), of
\(F-\E[F\mid W_{-j}]\). Each conditional variance equals half
the mean squared change on resampling \(W_j\).
By exchangeability this half equals the mean squared
positive change. Applying this to \(Q\) with the bounds
above gives \(\operatorname{Var}Q=O(\mu)\).
Fixed positive success at \(C\mu\) implies
\(\E Q\le C\mu+O(\sqrt\mu)\) by Chebyshev, then
\(Q\le C'\mu\) with probability tending to one.
The optimization can be performed directly on the observed
terminal counts.

\subsection*{Relays and individual pairs}
On the fixed-grid occupancy event each cell has between
\(c\mu\) and \(C\mu\) points of each color. Round each load
between its original floor and ceiling while preserving integer
divergence. For this, the subgraph of fractional edges has no
vertex of degree one, since the total divergence is integral
at every vertex. If this subgraph is nonempty it contains a
cycle. Change the loads by a signed constant along that cycle
until one hits an endpoint of its allowed interval, and repeat.
The face gauges increase by at most a constant. Remove directed cycles and decompose
the resulting integer flow into unit paths from surplus
to deficit cells. Reserve the corresponding red and blue
points at the endpoints. In cell \(v\), let
\(M_v=\min(R_v,B_v)\) be the remaining number of possible
red-blue relays. For internal visits entering in coordinate
direction \(i\), let \(T_{v,i}\) count them and give each
weight \(w_{v,i}=\min(1,M_v/(2T_{v,i}))\) when \(T_{v,i}>0\).

In the ordered list of internal visits of each path, group
successive visits until their weight first reaches one,
then repeat. Each completed group has weight less than two.
Choose one relay in each completed group. Fractionally this
is feasible by assigning a relay proportionally to the visit
weights, because the total weight offered by cell \(v\) is
at most \(M_v\). To make this choice integral, use a network
from groups of demand one to the eligible cells of capacity
\(M_v\). The fractional choice shows that all demand can
be met. With integer capacities, augmenting paths from
zero can use integer increments until demand is met,
by the residual cut argument above. A relay uses a blue point for the incoming
segment and a red point for the outgoing segment.
Pair the points along the resulting segments, and pair
all leftover points within their own cells.

Between successive relays or endpoints, the internal visits
have total weight at most four. The coordinate displacement
in units of \(h\) is at most two plus the number of internal
visits entering in that coordinate direction.
For \(\Psi(t)=\exp(t^{r_i})-1\) and \(A>0\), Jensen with denominator six
therefore charges the modular at displacement scale \(6hA\)
to a constant \(\Psi(1/A)\) and to
\(w_{v,i}\Psi(1/(Aw_{v,i}))\) over these visits.
Each cell is charged total weight \(O(\mu)\).
Moreover \(1/w_{v,i}\) is bounded by a constant times one
plus the sum of its incoming direction-\(i\) face magnitudes
divided by \(\mu\). Since \(\mu/N=h^2\), the Luxemburg
triangle inequality gives the point-coordinate norm bound
\[
 \|\Delta_i\|_{\psi_{r_i},\mathrm{points}}
 \le C h(1+\|q_i\|_{\psi_{r_i},\mathrm{faces}}/\mu).
\]
Here the load sum at a cell involves at most two adjacent
direction-\(i\) faces, and fixed modular prefactors are
absorbed by convexity under a fixed dilation. Both coordinates
refer to the same permutation. Since \(h\asymp s_N\) and the
optimized face norms are \(O(\mu)\) with high probability,
this proves Theorem~\ref{thm:finite}.

\begin{thebibliography}{9}
\bibitem{TalagrandProblems}
M.~Talagrand, \emph{My favorite problems}, manuscript, 2026.

\bibitem{umc:repository}
\url{https://github.com/grottofrancesco/UltimateMatchingTheorem}.
\end{thebibliography}

\end{document}
